\section{Correlações espúrias}

Modelos de aprendizado de máquina fundamentados na Minimização do Risco Empírico (ERM) são otimizados para reduzir a perda média global sobre o conjunto de treinamento. Entretanto, a presença de atributos espúrios desbalanceados frequentemente induz esses modelos a aprenderem ``atalhos'' preditivos em vez de características robustas. A Figura \ref{fig:spurious_correlation_sample} ilustra as consequências dessa dependência diante de uma mudança de distribuição: enquanto o conjunto de treino exibe uma forte correlação espúria entre a classe (pássaro) e o contexto (fundo), o conjunto de teste dissocia essa relação. Consequentemente, modelos que aparentam alta performance durante o treinamento falham ao generalizar para cenários onde a correlação espúria não se sustenta, resultando em uma degradação drástica de desempenho no pior grupo.

Para mitigar esse problema, abordagens como Group DRO (\textit{Distributionally Robust Optimization}) \cite{group_dro} buscam minimizar a perda do pior grupo em vez da perda média, geralmente através da reponderação de amostras de grupos minoritários. Contudo, a principal limitação dessas técnicas é a necessidade de anotações de grupo durante o treinamento, uma informação raramente disponível em cenários práticos.


\begin{figure}[htbp]
    \centering
    \includegraphics[width=0.8\textwidth]{images/spurious_problem.png}
    \caption{Ilustração do fenômeno de aprendizado de atalhos. Devido ao forte viés de seleção no conjunto de treinamento, o modelo pode utilizar o contexto de fundo como discriminador primário da classe, em detrimento das características intrínsecas do objeto. Isso resulta em falhas de predição quando a correlação espúria é quebrada no teste.}
    \label{fig:spurious_correlation_sample}
\end{figure}


\section{Aprendizado Contrastivo e SimCLR}

O Aprendizado Contrastivo é um paradigma de SSL onde dados não rotulados são contrastados uns com os outros. O objetivo é que representações de instâncias similares (pares positivos) sejam aproximadas no espaço latente, enquanto as de instâncias dissimilares (pares negativos) sejam afastadas.

O SimCLR \cite{simclr_main_paper} é um método proeminente que implementa essa abordagem. No SimCLR, cada amostra de um conjunto de dados não rotulado passa por uma série de transformações de aumento de dados, como recorte aleatório, espelhamento horizontal e vertical, desfoque Gaussiano e outras.

Uma rede neural convolucional, referida como o codificador $f(.)$, extrai vetores de características das imagens aumentadas. Para cada amostra original, um par positivo $(x_i, x_j)$ é criado usando os vetores de características de duas versões transformadas da mesma imagem (conforme ilustrado no Par 1 da Figura \ref{fig:simclr_pairs}). Pares negativos são formados a partir dos vetores de características de outros exemplos dentro do mini-lote, representando amostras dissimilares. Esses vetores de características são então passados por uma \textit{cabeça de projeção} $g(.)$, uma pequena rede neural que mapeia as representações extraídas para um espaço de dimensão inferior. Os vetores projetados resultantes, $z_i$ e $z_j$, são então usados na função de perda contrastiva.

\begin{equation*}
\ell_{i,j} = -\log \frac{\exp(\text{sim}(z_i, z_j)/\tau)}{\sum_{k=1}^{2N} \mathbb{1}_{[k \neq i]} \exp(\text{sim}(z_i, z_k)/\tau)}
\end{equation*}

A função indicadora $\mathbb{1}_{[k \neq i]}$ assume o valor 1 quando $k \neq i$, garantindo que apenas amostras negativas contribuam para o denominador. O hiperparâmetro $\tau$ atua como um fator de temperatura, regulando a escala das pontuações de similaridade. A perda total é computada pela média das contribuições de todos os pares positivos no mini-lote, considerando ambas as ordenações de pares possíveis $(i, j)$ e $(j, i)$.

\begin{figure}[h]
    \centering
\includegraphics[width=0.8\textwidth]{images/simclr_pairs.png}
    \caption{Ilustração de pares positivos e negativos no \textit{framework} SimCLR. O \textbf{Par 1 (positivo)} mostra duas aumentações da mesma imagem. Os \textbf{Pares 2, 3 e 4 (negativos)} mostram a imagem comparada com outras amostras distintas do conjunto de dados.}
    \label{fig:simclr_pairs}
\end{figure}


\section{Inteligência Artificial Explicável e Grad-CAM}

Técnicas de Inteligência Artificial Explicável \cite{xai_survey} podem ajudar a identificar se os modelos aprenderam algum viés indesejado. Essas técnicas podem fornecer explicações em nível global ou local. A explicação global visa entender o modelo como um todo, revelando as regras por trás de cada decisão. Por outro lado, a explicação local foca em identificar quais características de uma amostra específica foram mais importantes para gerar uma determinada classificação. Entre os métodos locais de XAI, o Grad-CAM \cite{gradcam_main_paper} se destaca.

O Grad-CAM é uma técnica derivada do \textit{Class Activation Mapping} (CAM) \cite{cam_main_paper}, cuja principal diferença é poder ser aplicada a qualquer arquitetura de \textit{Convolutional Neural Network} (CNN) \cite{cnn_introduction} sem exigir modificações na mesma, o que a torna amplamente utilizada. A técnica calcula o gradiente ($\frac{\partial y^c}{\partial A^k}$) da pontuação $y^c$ para a classe $c$ em relação à ativação do mapa de características $A^k$ da $k$-ésima camada convolucional. Esses gradientes são então agregados globalmente através de \textit{max-pooling}, permitindo determinar a importância de cada convolução para a classe, representada como $\alpha_c^k$. A saída resultante é um mapa de localização $L^c$, obtido pela aplicação de uma soma ponderada seguida por uma função de ativação ReLU, que enfatiza apenas as contribuições positivas (visualizado na Figura \ref{fig:gradcam_sample}).

\begin{figure}[h]
    \centering
\includegraphics[width=0.8\textwidth]{images/gradcam_sample.png}
    \caption{Visualização das explicações via Grad-CAM. A imagem mostra a entrada original, o mapa de calor gerado pelos gradientes da última camada convolucional e a sobreposição final, destacando as regiões que mais contribuíram para a predição da classe.}
    \label{fig:gradcam_sample}
\end{figure}

\section{Explanation-Guided Learning}

Tradicionalmente, técnicas de XAI são post-hoc (aplicadas após o treino), o que pode exigir múltiplas iterações para identificar e mitigar vieses. Contudo, uma nova linha de pesquisa vem ganhando força: \textit{Explanation-Guided Learning} (ou \textit{Model Guidance}). Essa abordagem propõe regularizar o modelo durante o próprio treinamento, utilizando explicações de um mecanismo XAI.

Em \cite{gradmask}, por exemplo, os autores introduziram o GradMask, uma técnica de EGL que penaliza mapas de saliência, derivados dos gradientes, quando são inconsistentes com áreas de interesse na imagem. O objetivo é mitigar a superestimação de características não relacionadas ao tumor na classificação de imagens médicas — um desafio comum em conjuntos de dados pequenos ou quando artefatos hospitalares afetam as previsões. Os autores relataram que o GradMask melhora a acurácia no teste em 1-3\% em comparação ao modelo base, reduzindo o \textit{overfitting} e aumentando a robustez do modelo a variações nos dados.

Uma das abordagens seminais nesta área é o ``\textit{Right for the Right Reasons}'' (RRR) \cite{rrr}. A técnica propõe regularizar o modelo de forma eficiente, examinando e penalizando seletivamente os gradientes de entrada. Na prática, um especialista pode fornecer uma máscara binária, indicando quais dimensões da entrada são irrelevantes para uma determinada amostra. Um termo de penalização é então adicionado à função de perda principal, forçando a magnitude desses gradientes a ser próxima de zero nas regiões marcadas como irrelevantes. Isso força o classificador a aprender explicações alternativas e a basear suas decisões nas características corretas, melhorando a generalização em cenários com mudança de distribuição.

Seguindo uma linha similar, \cite{gals} propõem o GALS (\textit{Guiding visual Attention with Language Specification}). A abordagem utiliza especificações de linguagem de alto nível (por exemplo, ``uma foto de um pássaro'') e um modelo de visão-linguagem pré-treinado (como o CLIP \cite{clip}) para gerar mapas de atenção que destacam as regiões relevantes. Esses mapas servem então como supervisão para o classificador principal durante o treinamento. Ao adaptar a técnica RRR, o GALS adiciona uma perda de atenção auxiliar que guia a atenção espacial do classificador para longe dos distratores. O método demonstrou melhorias significativas na acurácia do pior grupo e em métricas de \textit{fairness} em conjuntos de dados com viés explícito e implícito.

Mais recentemente, em \cite{crayon}, os autores introduziram o CRAYON, uma técnica inovadora para corrigir a atenção do modelo usando simples anotações de ``sim'' ou ``não''. Ao contrário dos métodos tradicionais, que exigem rótulos complexos pixel a pixel ou anotações explícitas para áreas irrelevantes, o CRAYON oferece uma solução escalável e prática. Ele aprimora tanto técnicas de interpretação clássicas quanto modernas, como o CRAYON-Attention, que refina os mapas de saliência para focar em regiões relevantes, e o CRAYON-Pruning, que remove neurônios irrelevantes. Um extenso conjunto de experimentos, incluindo avaliações quantitativas e humanas, demonstrou o desempenho superior do CRAYON, superando 12 métodos existentes em três conjuntos de dados de benchmark e provando ser mais eficaz do que métodos que exigem anotações mais complexas.

\section{Cross-Modal Learning}

O aprendizado \textit{Cross-Modal} (ou multimodal) é um paradigma de aprendizado de máquina que treina modelos utilizando simultaneamente diferentes modalidades de dados, como imagens e texto, visando melhorar o desempenho em tarefas específicas. Uma das hipóteses centrais deste trabalho de doutorado é que certas modalidades de dados são inerentemente menos suscetíveis a correlações espúrias. Por exemplo, enquanto um modelo de visão pode ser enviesado por características de fundo, as descrições textuais correspondentes geralmente focam no objeto de interesse, oferecendo um sinal de supervisão mais limpo.

Essa abordagem tem demonstrado resultados robustos na literatura, sendo o CLIP (\textit{Contrastive Language-Image Pre-training}) \cite{clip} uma de suas implementações mais proeminentes. O CLIP aprende representações visuais transferíveis a partir de supervisão em linguagem natural, sendo pré-treinado em uma escala massiva de 400 milhões de pares de imagem-legenda coletados da web. Sua arquitetura consiste em dois \textit{encoders} distintos, um para imagens e outro para texto, que mapeiam ambas as entradas para um espaço de \textit{embedding} compartilhado. Um objetivo de aprendizado contrastivo é utilizado para garantir que os \textit{embeddings} de pares imagem-texto correspondentes sejam aproximados nesse espaço, enquanto pares não correspondentes sejam afastados.

Um princípio semelhante é explorado por \cite{hager2023best}, que introduz uma arquitetura de aprendizado contrastivo multimodal combinando imagem e dados tabulares. A premissa é aproveitar dados multimodais de biobancos para pré-treinar \textit{encoders} que melhorem a predição unimodal (feita apenas com imagens) em cenários clínicos onde dados tabulares extensos são inviáveis. O método utiliza um \textit{framework} de aprendizado contrastivo auto-supervisionado, baseado em SimCLR e SCARF. O estudo revela que características tabulares morfométricas, como o volume ventricular e outras que descrevem tamanho e forma, são desproporcionalmente importantes. A razão é que elas possuem ``correlatos diretos'' nas imagens, o que facilita a minimização da perda contrastiva e melhora a qualidade das representações aprendidas.

A versatilidade das abordagens contrastivas multimodais estende-se com sucesso a outros domínios, como a química. Um exemplo notável é o framework SMICLR \cite{smiclr}, que propõe o treinamento conjunto de dois encoders distintos para processar a mesma molécula: um baseado em grafos moleculares e outro em sequências de texto SMILES. Ao alinhar essas modalidades complementares via objetivo contrastivo, o modelo captura representações mais robustas das relações estrutura-propriedade. A eficácia do método foi comprovada no dataset QM9, onde o fine-tuning resultou em uma redução de erro médio de 44\% para propriedades energéticas e 25\% para eletrônicas. Esses resultados, superiores aos baselines puramente supervisionados, validam o potencial do aprendizado multimodal na extração de características latentes complexas em cenários científicos.

\section{Abordagens mais diretas para lidar com corelações espúrias}

Embora estratégias como \textit{Explanation-Guided Learning} e aprendizado \textit{Cross-Modal} sejam abordagens viáveis para mitigar correlações espúrias, existem paradigmas alternativos focados na robustez. Um dos métodos mais proeminentes é o Group DRO (\textit{Distributionally Robust Optimization}) \cite{group_dro}, que redefine o objetivo de treinamento: em vez de minimizar a perda média (ERM), ele busca minimizar a perda do pior caso em um conjunto de grupos pré-definidos.

Esses grupos são explicitamente formados usando conhecimento prévio sobre potenciais vieses, geralmente combinando o rótulo da tarefa com um atributo espúrio. O método foca a otimização em reduzir a perda máxima identificada entre esses grupos durante o treinamento. Contudo, em redes neurais sobreparametrizadas que podem zerar facilmente a perda de treino, o \textit{Group DRO} só se torna eficaz quando acoplado a uma regularização forte, como uma penalidade $l_2$ elevada ou \textit{early stopping}. Essa regularização impede que o modelo se ajuste perfeitamente aos dados, forçando o otimizador DRO a fazer concessões. Isso garante um bom desempenho no pior grupo de treino, o que se traduz em melhor generalização para esse grupo no teste.

Abordando a principal limitação do Group DRO — a necessidade de rótulos de grupo pré-definidos — uma outra estratégia proeminente é o Correct-N-Contrast (CNC) \cite{cnc}. Este método opera no cenário mais desafiador onde os rótulos de atributos espúrios não estão disponíveis durante o treinamento. O CNC é uma abordagem de dois estágios  que utiliza o próprio viés dos modelos padrão a seu favor. No primeiro estágio, um modelo ERM é treinado, partindo da observação de que ele aprende as correlações espúrias e, assim, torna-se um bom preditor desses atributos espúrios. No segundo estágio, as previsões desse modelo ERM são usadas como ``pseudo-rótulos'' para guiar um aprendizado contrastivo. O objetivo é forçar o novo modelo a ``puxar para perto'' as representações de amostras que possuem a mesma classe, mas que o modelo ERM previu de forma diferente (indicando atributos espúrios distintos). Simultaneamente, ele ``empurra para longe'' amostras de classes diferentes que o modelo ERM previu da mesma forma (indicando atributos espúrios similares), aprendendo assim representações robustas.

Mais recentemente \cite{memorization}, uma linha de pesquisa alternativa passou a investigar o problema sob a ótica da memorização. Esta abordagem parte da observação de que o desempenho fraco em grupos minoritários está ligado à ``memorização espúria'': a tendência das redes neurais de memorizar exemplos atípicos (minoritários) usando apenas um subconjunto muito pequeno de ``neurônios críticos''. Embora essa memorização localizada garanta alta acurácia de treino para esses grupos, ela falha em generalizar, levando a um desempenho de teste ruim. Para intervir nesse mecanismo, os autores propuseram um \textit{framework} de \textit{fine-tuning} que utiliza um ``modelo alvo'' e um ``modelo auxiliar''. O modelo auxiliar é uma versão podada do modelo alvo, onde esses neurônios críticos (identificados por alta magnitude ou gradiente) são mascarados. Utilizando um objetivo de aprendizado contrastivo, o \textit{framework} força o modelo alvo a alinhar suas representações de características com as do modelo auxiliar, que é incapaz de usar a memorização espúria. Ao perturbar ativamente essa memorização, o método obriga o modelo a aprender representações mais robustas para os grupos minoritários, melhorando significativamente a acurácia do pior grupo no teste.